% ECONOMICS_PROBLEM: {"id": "L94-74", "title": "Electricity-market design with high renewables and storage", "jel_code": "L94", "source_id": "OP-0115"}
% !TeX program = xelatex
\documentclass[11pt,a4paper]{article}
\usepackage[margin=23mm]{geometry}
\usepackage{fontspec}
\setmainfont{Latin Modern Roman}
\usepackage{amsmath,amssymb,mathtools}
\usepackage{xcolor,enumitem,booktabs,longtable,array,xurl,hyperref}
\definecolor{muted}{HTML}{53636C}
\hypersetup{colorlinks=true,urlcolor=blue,linkcolor=blue}
\setlength{\parindent}{0pt}
\setlength{\parskip}{5pt}
\newcommand{\R}{\mathbb R}
\newcommand{\E}{\mathbb E}
\newcommand{\N}{\mathbb N}
\newcommand{\ind}{\mathbf 1}
\newcommand{\Var}{\operatorname{Var}}
\newcommand{\Cov}{\operatorname{Cov}}
\newcommand{\supp}{\operatorname{supp}}
\DeclareMathOperator*{\argmin}{arg\,min}
\DeclareMathOperator*{\argmax}{arg\,max}
\newcommand{\entrymeta}[3]{{\sffamily\small\color{muted}\textbf{\texttt{#1}}\quad|\quad #2\par #3\par}}
\newcommand{\Problemref}[1]{\texttt{#1}}
\newcommand{\JELref}[2]{\texttt{#2} (JEL \texttt{#1})}
\begin{document}
\section*{Electricity-market design with high renewables and storage}
\textbf{Problem ID:} \texttt{L94-74}\quad \textbf{JEL:} \texttt{L94}\par
% BEGIN ECONOMICS STATEMENT
\label{problem:OP-0115}
\label{audited:ECON-040}


\label{atlas:prob:C5}

\noindent\textbf{Original identifier:} C5 (atlas item 115). \textbf{Field:} Environmental and climate economics. \textbf{Classification:} editorial research family with established restricted solutions; current open status not certified.

\subsection*{Definitions and assumptions}

Take locations $\ell=1,\ldots,L$, periods $t=1,\ldots,H$ of duration $\Delta t$ hours, and a stochastic renewable/demand process. At each location let conventional generation $g_{\ell t}$, renewable output $r_{\ell t}$, charging $c_{\ell t}$, discharging $d_{\ell t}$, demand $q_{\ell t}$ and unserved demand $u_{\ell t}$ be energies in MWh. Network transfers obey declared flow constraints and losses; nodal balance equates supply plus net imports to $q_{\ell t}-u_{\ell t}+c_{\ell t}$. Storage energy $b_{\ell t}$ obeys $b_{\ell,t+1}=b_{\ell t}+\eta_c c_{\ell t}-d_{\ell t}/\eta_d$, with efficiencies in $(0,1]$, energy-capacity bounds and power limits multiplied by $\Delta t$. Commitment, ramping, investment and degradation costs are specified in a fixed real-currency year.

\subsection*{Formal research question}

Characterize market rules implementing reliable, cost-effective investment and dispatch with strategic generators, storage owners and flexible demand. Let $R$ denote a complete rule system, including price formation, capacity payments and settlement, and $\mathcal E(R,P)$ its equilibria under scenario law $P$. For a stated ambiguity class $\mathcal P$, reliability bound $\varepsilon$ in expected unserved MWh over the horizon, and total social cost $C(e,P)$ including emissions and lost load, seek
\[
\inf_R\ \sup_{P\in\mathcal P}\sup_{e\in\mathcal E(R,P)} C(e,P)
\quad\text{subject to}\quad
\sup_{P\in\mathcal P,e\in\mathcal E(R,P)}
\mathbb E_P\!\left[\sum_{\ell,t}u_{\ell t}\right]\le\varepsilon .
\]
State participation and budget restrictions, and either prove equilibrium existence or exclude rule systems lacking equilibrium. Compare market outcomes to a planner using the same physical and information constraints.

\subsection*{Interpretation and scope}

Zero fuel cost during some hours does not eliminate scarcity, network congestion, opportunity costs of stored energy or financing costs. In a convex price-taking benchmark, nodal and intertemporal prices can decentralize a planner; binary unit commitment and strategic behavior require additional analysis. Reliability metrics such as expected unserved energy and loss-of-load probability are not interchangeable. Emissions costs must use a declared price per tonne rather than mix tonnes of carbon with tonnes of $\mathrm{CO}_2$. Estimated storage capacity is evidence about technology availability, not a welfare or equilibrium theorem. Demand response creates a resource only after discomfort, privacy and adjustment costs are specified.

\subsection*{Source audit and status}

[S1] surveys restructuring, [S2] quantifies intermittency in an electricity model, and [S3] discusses wholesale-market challenges at renewable scale. These establish substantial model-specific analyses. The input's JRC inventory link supports storage data, with some MWh capacities estimated; it does not solve design. The robust strategic-market question above is editorial, and its exact 2026 status is not certified by these older sources.

\subsection*{References}

\begin{enumerate}[label={[\textnormal{S}\arabic*]},leftmargin=*]

\item Severin Borenstein and James Bushnell (2015). \emph{The US Electricity Industry After 20 Years of Restructuring}. Annual Review of Economics 7, 437–463. \url{https://doi.org/10.1146/annurev-economics-080614-115630}. \textbf{Locator:} Publisher or institutional abstract and bibliographic record. \textbf{Support and limits:} Restructuring survey; does not establish a universally optimal renewables/storage design.

\item Gautam Gowrisankaran, Stanley S. Reynolds and Mario Samano (2016). \emph{Intermittency and the Value of Renewable Energy}. Journal of Political Economy 124(4), 1187–1234. \url{https://doi.org/10.1086/686733}. \textbf{Locator:} Publisher or institutional abstract and bibliographic record. \textbf{Support and limits:} Intermittency costs in an estimated electricity model; low marginal generation cost does not remove scarcity rents.

\item Paul L. Joskow (2019). \emph{Challenges for wholesale electricity markets with intermittent renewable generation at scale: the US experience}. Oxford Review of Economic Policy 35(2), 291–331. \url{https://doi.org/10.1093/oxrep/grz001}. \textbf{Locator:} Publisher or institutional abstract and bibliographic record. \textbf{Support and limits:} Wholesale-market design challenges; reliability and investment require explicit institutional assumptions.

\end{enumerate}

\subsection*{Integrated refinement: ECON-040}
{\sffamily\small\color{muted}Spatial energy networks, storage, and the green transition\par}
This refinement adopts the additional source conventions
(page~\pageref{audited:conventions}), subject to its local restrictions.
\textbf{Spatial investment and storage benchmark for the market rules.}
Specify edge capacities, delivered flow fractions and the chosen AC/DC
power-flow model; a generic flow graph alone does not define an electricity
system. Storage at node $r$ obeys
\[
b_{r,t+1}=(1-\rho_r)b_{rt}+\eta_r^{\rm in}q_{rt}^{\rm in}
                              -q_{rt}^{\rm out}/\eta_r^{\rm out},
\]
with bounded inventories and charging/discharging, efficiencies in $(0,1]$,
and declared initial and terminal stocks. For feasible investment
$k\in\mathcal K$ and nonanticipative operation, characterize
$\inf_{k\in\mathcal K}C(k,\tau)$ and regional welfare, including investment,
shortage, emissions and equilibrium economic responses. Use this planner as the
benchmark for decentralization by the existing market rules. Specify investment
versus dispatch information dates, weather, endogenous demand and location,
physical flow/loss restrictions, and attainment conditions. Expected shortage
penalties and reliability constraints define different comparisons.
\subsection*{Supplied source audit for ECON-040}
Local [S1], [S2], etc. in this audit and the references immediately below
belong to ECON-040, independently of the original entry's reference numbering.
[S1], Section 7.2, printed p. 298, calls for energy trade, storage and transmission. The editorial constraints distinguish electrical physics from generic transport flow.
\subsection*{References for integrated refinement ECON-040}
{\footnotesize\raggedright
\par\noindent\textbf{[S1]} Klaus Desmet and Esteban Rossi-Hansberg (2024). \emph{Climate Change Economics over Time and Space}. Annual Review of Economics 16, 271–304. DOI: 10.1146/annurev-economics-072123-044449.\par \textit{Verified locator / source role:} Section 7.2, Green Transition and Trade in Energy, p. 298.\par \url{https://rossihansberg.economics.uchicago.edu/CCETS.pdf}\par\smallskip
}

\subsection*{Common conventions: Source mathematical conventions}

Unless an entry states otherwise, agent and firm sets are finite. State and action spaces are measurable, strategies and outcome maps are measurable, probability laws are specified, and expectations exist for the targets under discussion. Infinite-horizon discount factors lie in $(0,1)$ and discounted objectives are finite. Norms, tolerances and complexity input sizes must be specified for computational comparisons. Constants or primitives that appear locally are inputs, not empirical facts asserted by this catalogue.

\textbf{Identification.} A statistical model $m\in\mathcal M$ implies an observable law $P_m$ and target $\theta(m)$. At law $P$, its identified set is
\[
\Theta(P;\mathcal M)=\{\theta(m):m\in\mathcal M,\ P_m=P\}.
\]
Point identification holds when this set is a singleton, or equivalently when $P_{m_1}=P_{m_2}$ implies $\theta(m_1)=\theta(m_2)$ for all compatible models. Sharp bounds mean bounds attained, or approached when the boundary is not attained, by compatible models. Writing this set is a definition, not a solution: the substantive task is to establish informative restrictions and derive or estimate the set. An unrestricted-confounding model may give only trivial bounds.

\textbf{Inference.} Identification, estimability and valid uncertainty quantification are separate questions. Fix $\alpha\in(0,1)$. A confidence procedure $C_n$ over a declared law class $\mathcal P$ has asymptotic uniform coverage at level $1-\alpha$ when
\[
\liminf_{n\to\infty}\inf_{P\in\mathcal P}\Pr_P\{\theta(P)\in C_n\}\geq1-\alpha.
\]
For set-identified targets the coverage event must be defined separately, for example coverage of the full identified set. If an entry uses identification-indexed models instead of uniquely indexed laws, the same coverage convention is applied over those models. Pointwise limits do not establish uniform validity, and asymptotic validity does not guarantee finite-sample precision. Sample sizes, dependence structures and weak-identification sequences are part of each application.

\textbf{Counterfactuals.} $Y_i(a)$ is a potential outcome under a specified intervention, including its timing, implementation and versions. It is not $Y_i$ conditional on voluntarily choosing $a$. A full intervention vector permits interference; shorthand scalar treatment notation does not silently impose no interference. Assignment, selection, attrition, anticipation and measurement must be specified. A claim about an instrument's local effect is not automatically a population policy effect.

\textbf{Economic equilibrium.} A counterfactual model includes best responses, constraints, feasibility, market clearing, state transitions and expectations consistent with the stated information. When several equilibria exist, either a declared selector is an input or the target is set-valued. Comparisons across policy regimes must use the stated selection convention. Reduced-form relationships are not presumed invariant after an equilibrium-changing intervention.

\textbf{Optimization and welfare.} Optimization is over the stated feasible set. If attainment is not established, $\sup$ or $\inf$ and $\varepsilon$-optimal choices replace an asserted maximizer or minimizer. For example, with value $V=\sup_{a\in\mathcal A}W(a)<\infty$, an $\varepsilon$-optimal set is $\{a\in\mathcal A:W(a)\geq V-\varepsilon\}$ for $\varepsilon>0$. Benefits, costs, risk and health or environmental outcomes can be aggregated only using declared units and conversion weights. Interpersonal utility weights, the treatment of fiscal transfers and foreign profits, discounting, and the behavioral welfare criterion are explicit inputs. A comparative-static derivative additionally requires existence and appropriate differentiability of the selected counterfactual.

\textbf{Sources.} Local labels [S1], [S2], and so forth restart in every question. Locators identify the relevant abstract, model, section or result at the level actually checked; a bibliographic or abstract-level verification is not represented as inspection of an entire proof. Working papers are distinguished from later publications and changing versions. Source summaries are paraphrased. All URL checks and editorial formulations refer to the audit date stated above.

\subsection*{Common conventions: Additional-source status and mathematical conventions}

\label{audited:conventions}

Of the 100 supplied ECON entries, 65 distinct problems appear here; 32 close
matches refine earlier entries, and three duplicate questions map to existing
entries. Appendix~\ref{app:audited-crosswalk} lists every destination.
The attachment's P/R/S/U distinctions and source audits are retained. Its
precise mathematical formalizations are editorial unless the individual audit
says otherwise. Candidate subsidy targets ECON-004--006 retain their U flags;
the resolved approval-core question ECON-007 updates OP-0202 above.
The following supplied conventions govern both this part and the attributed
ECON refinements elsewhere in the catalogue, subject to local restrictions.

\textbf{Parameterized questions.} A problem family lists inputs that must be fixed before an instance is evaluated: population, horizon, policies, information, data law, model class and welfare weights. These are required choices, not quantities the citations are assumed to specify. Undefined applications should not be treated as identified estimands.

\textbf{Indices and integrability.} Sets of agents, firms and regions are finite unless a population measure is explicitly used. \(i,j\) index units, \(r\) regions, \(t\) dates and \(h\) horizons. \(\mathbb E\) and \(\Pr\) refer to the declared population and law; \(\mathbf1\{A\}\) is an indicator. Every displayed expectation and discounted sum needs existence in the stated domain. Infinite horizons use a discount in \((0,1)\) and a convergence condition; finite horizons may use discount one.

\textbf{Optimization and existence.} A displayed \(\max\), \(\min\), or \(\arg\max\) is conditional on attainment. For finite feasible menus this is immediate; general spaces need continuity and compactness/coercivity or another existence argument. Without attainment, the target is the supremum/infimum and arbitrarily near-optimal feasible choices. \(\inf\varnothing=+\infty\). Empty feasibility and an unidentified target are different issues.

\textbf{Potential outcomes.} \(Y_i(z)\) is the outcome under a specified intervention, not the observed conditional mean among people choosing \(z\). In binary comparisons \(\Delta Y_i=Y_i(1)-Y_i(0)\). Specify assignment, treatment versions, compliance, information and observation horizon. Interference is allowed under a full policy vector; compressed exposure notation needs a justified mapping. No interference, consistency and overlap are substantive assumptions, not automatic consequences of notation.

\textbf{Populations and observation.} Fix baseline eligibility and population weights. Conditioning on post-policy employment, survival, relocation or program uptake can change the population being compared. Death is not ordinary loss to follow-up; outcomes truncated by death need a composite convention, joint survival target, or explicitly defined survivor stratum. Fertility-changing interventions need a population functional rather than an unexplained fixed descendant cohort. Measurement and missingness models must be supplied.

\textbf{Identification.} A structural model \(m\in\mathcal M\) implies observable law \(P_m\) and target \(\theta(m)\). Identification means
\[
 P_{m_1}=P_{m_2}\ \Longrightarrow\ \theta(m_1)=\theta(m_2)
 \quad\text{for all }m_1,m_2\in\mathcal M .
\]
Without identification, the target is
\[
 \Theta(P;\mathcal M)=\{\theta(m):m\in\mathcal M,\ P_m=P\}.
\]
Writing \(\theta(P)\) before establishing identification can hide incompatible latent models. Confidence sets additionally need a sampling law, confidence level, asymptotic sequence and uniformity class. Point identification, coverage and informative precision are separate properties.

\textbf{Mediation.} Total effects of randomized assignment are distinct from cross-world natural indirect effects. Belief/effort-channel effects require a meaningful mediator intervention and additional measurement, exclusion or mediation assumptions. Randomization of the initial policy alone does not supply those assumptions. Report bounds or interventional alternatives when the stronger target is unsupported.

\textbf{Equilibrium.} A counterfactual \(W(z)\), \(p(z)\), or \(\mu_t^z\) requires individual optimization, feasibility, government/resource budgets, market clearing and state-transition laws. State equilibrium existence and selection, or report ranges over compatible equilibria. Initial-state integration includes subsequent endogenous transitions; current-state integration must use the policy-dependent distribution. Reduced-form invariance after policy is not assumed.

\textbf{Welfare accounts.} Scores, utility, health, dollars and ecological quantities require explicit conversion before addition. Fees, taxes, subsidies, wages and extortion payments are transfers between parties; distributional weights can make them welfare relevant, but their face value is not automatically a resource loss. State ownership, geographic boundaries, foreign-income weights and fiscal finance. Do not add profits to household welfare already including dividends, or count land capitalization and the underlying benefit twice. A benefit-cost expression must distinguish gross benefits from net welfare and subtract every real cost exactly once.

\textbf{Derivatives and risk.} Log derivatives require positive arguments; local responses require admissible perturbations and specified equilibrium branches. Differentiation under expectations needs regularity/domination, and boundaries may allow only one-sided derivatives. Risk uses a specified law; ambiguity uses a declared nonempty law set on a common history space. Adaptive policies must be nonanticipative. A robust criterion is an editorial choice unless attributed explicitly.

\textbf{Scope overlaps.} Allocation, collective choice, contracts, household bargaining and coordinated policy overlap economics with optimization and game theory. They are retained because they appeared in the original catalogue; no claim of a strictly game-theory-free selection is made.
% END ECONOMICS STATEMENT
\end{document}
